% Part: first-order-logic
% Chapter: tableaux
% Section: derivations

\documentclass[../../../include/open-logic-section]{subfiles}

\begin{document}

\iftag{FOL}
{\olfileid[id]{fol}{tab}{der}}
{\olfileid[id]{pl}{tab}{der}}

\olsection{\usetoken{P}{tableau}}

\begin{explain}
Kita telah menjelaskan apa itu asumsi dan memberikan aturan-aturan
inferensi. !!^{tableau}s dibangkitkan secara induktif dari keduanya:
setiap !!{tableau} merupakan satu cabang yang terdiri atas satu atau
lebih asumsi, atau dihasilkan dari !!a{tableau} dengan menerapkan salah
satu aturan inferensi pada suatu cabang.
\end{explain}

\begin{defn}[!!^{tableau}]
!!^a{tableau} untuk asumsi-asumsi $\sFmla{S_1}{!A_1}$, \dots,
$\sFmla{S_n}{!A_n}$ (dengan setiap $S_i$ berupa $\True$ atau~$\False$)
adalah pohon berhingga yang terdiri atas !!{signed formula}s dan memenuhi
syarat-syarat berikut:
\begin{enumerate}
\item Sebanyak $n$ !!{signed formula}s paling atas pada pohon adalah
  $\sFmla{S_i}{!A_i}$, disusun satu di bawah yang lain.
\item Setiap !!{signed formula} dalam pohon yang bukan salah satu asumsi
  dihasilkan melalui penerapan yang benar dari suatu aturan inferensi
  pada !!a{signed formula} di cabang di atasnya.
\end{enumerate}
Suatu cabang dari !!a{tableau} \emph{tertutup} jika dan hanya jika
cabang tersebut memuat $\sFmla{\True}{!A}$ dan~$\sFmla{\False}{!A}$,
dan \emph{terbuka} jika tidak demikian. !!^a{tableau} yang setiap
cabangnya tertutup adalah \emph{!!{tableau} tertutup} (untuk himpunan
asumsinya). Jika !!a{tableau} tidak tertutup, yaitu jika memuat
sekurang-kurangnya satu cabang terbuka, tableau tersebut \emph{terbuka}.
\end{defn}

\begin{ex}
  Setiap himpunan asumsi dengan sendirinya merupakan !!a{tableau}, tetapi
  pada umumnya tidak tertutup. (Jelas bahwa tableau itu tertutup hanya
  jika asumsi-asumsinya sudah memuat sepasang !!{signed formula}s
  $\sFmla{\True}{!A}$ dan~$\sFmla{\False}{!A}$.)

  Dari !!a{tableau} (terbuka ataupun tertutup), kita dapat memperoleh
  tableau baru yang lebih besar dengan menerapkan salah satu aturan
  inferensi pada !!a{signed formula}~$\sFmla{S}{!A}$ di dalamnya. Aturan itu akan
  menambahkan satu atau lebih !!{signed formula}s ke ujung setiap cabang
  yang memuat kemunculan~$\sFmla{S}{!A}$ yang dikenai aturan tersebut.

  Sebagai contoh, perhatikan asumsi $\sFmla{\True}{!A \land \lnot
    !A}$. Berikut adalah !!{tableau} (terbuka) yang hanya terdiri atas
  asumsi tersebut:
  \begin{center}
    \begin{tableau}{}
      [\sFmla{\True}{\formula{A} \land \lnot \formula{A}}, just=\TAss]
    \end{tableau}{}
  \end{center}
  Kita memperoleh !!{tableau} baru darinya dengan menerapkan aturan
  $\TRule{\True}{\land}$ pada asumsi tersebut. Aturan itu memungkinkan
  kita menambahkan dua baris baru pada !!{tableau}, yaitu
  $\sFmla{\True}{!A}$ dan $\sFmla{\True}{\lnot !A}$:
  \begin{center}
    \begin{tableau}{}
      [\sFmla{\True}{\formula{A} \land \lnot \formula{A}}, just=\TAss
        [\sFmla{\True}{\formula{A}}, just={\TRule{\True}{\land}[1]},
          [\sFmla{\True}{\lnot \formula{A}}, just={\TRule{\True}{\land}[1]}
          ]
        ]
      ]
    \end{tableau}{}
  \end{center}
  Ketika menuliskan !!{tableau}s, kita mencatat aturan yang telah
  diterapkan di sebelah kanan (misalnya, $\TRule{\True}{\land} 1$
  berarti bahwa !!{signed formula} pada baris itu dihasilkan dengan
  menerapkan aturan $\TRule{\True}{\land}$ pada !!{signed formula} di
  baris~$1$). Hasilnya adalah !!{tableau} baru yang kini memuat
  !!{signed formula}s tambahan, tetapi aturan hanya dapat diterapkan pada satu di antaranya,
  yaitu $\sFmla{\True}{\lnot !A}$ (dalam hal ini, aturan
  $\TRule{\True}{\lnot}$). Penerapan itu menghasilkan !!{tableau}
  tertutup berikut:
  \begin{center}
    \begin{tableau}{}
      [\sFmla{\True}{\formula{A} \land \lnot \formula{A}}, just=\TAss
        [\sFmla{\True}{\formula{A}}, just={\TRule{\True}{\land}[1]}
          [\sFmla{\True}{\lnot \formula{A}}, just={\TRule{\True}{\land}[1]}
            [\sFmla{\False}{\formula{A}}, just={\TRule{\True}{\lnot}[3]}, close]
          ]
        ]
      ]
    \end{tableau}{}
  \end{center}
\end{ex}

\end{document}
